\documentclass[10pt,a4paper]{article}

\usepackage[utf8]{inputenc}
\usepackage[T1]{fontenc}
\usepackage[spanish]{babel}
\usepackage{geometry}
\geometry{margin=2cm, bottom=2.2cm}
\usepackage{parskip}
\usepackage{color}
\usepackage{graphicx}
\usepackage{fancyhdr}
\usepackage{lastpage}
\usepackage{titlesec}
\usepackage{tabularx}
\usepackage{booktabs}
\usepackage{hyperref}

% Minimalismo asertivo
\definecolor{gray75}{gray}{0.6}
\pagestyle{fancy}
\fancyhf{}
\fancyhead[L]{\small SSI -- Silent Speech Interface}
\fancyhead[R]{\small Plataforma de Desarrollo}
\fancyfoot[C]{\small \thepage /\pageref{LastPage}}
\renewcommand{\headrulewidth}{0.4pt}
\renewcommand{\footrulewidth}{0pt}

\titlespacing*{\section}{0pt}{1.2ex}{0.8ex}
\titlespacing*{\subsection}{0pt}{1ex}{0.5ex}
\titleformat{\section}{\large\bfseries}{\thesection}{1em}{}
\titleformat{\subsection}{\normalsize\bfseries}{\thesubsection}{1em}{}

\usepackage[scaled=0.92]{beramono}
\usepackage{listings}
\lstdefinestyle{mini}{basicstyle=\ttfamily\footnotesize,frame=none,
  numbers=none,breaklines=true}

\begin{document}

\section*{SSI -- Silent Speech Interface\\\smallskip
\small Plataforma de captura sincronizada para interfaz de habla silenciosa}

\vspace{-0.3cm}
\noindent\rule{\textwidth}{0.4pt}
\vspace{0.3cm}

\subsection*{Objetivo}

Construir un prototipo funcional de captura sincronizada de vibraciones laríngeas
y audio de referencia, para generar un dataset de pares
\texttt{(voz\_mic, vibracion\_acc)} que alimente un modelo de deep learning
embebido. El producto final es un dispositivo autónomo, portátil y discreto que
reconstruye la voz a partir de las vibraciones mínimas de la laringe.

\subsection*{Pipeline de Desarrollo}

\begin{lstlisting}[style=mini]
IIS3DWB ---SPI DMA (26.7 kHz)---+
                                 +-- NUCLEO-H723ZG --USB CDC--> PC --> .ssr --> .npz --> Colab
INMP441 ---I2S DMA (16 kHz)------+
\end{lstlisting}

\vspace{-0.3cm}
\section{Arquitectura del sistema}

\subsection*{Sensores}

\begin{tabularx}{\textwidth}{l X}
\toprule
\textbf{Componente} & \textbf{Rol} \\
\midrule
IIS3DWB & Acelerómetro triaxial SPI. ODR fijo 26.7\,kHz, BW 6.3\,kHz.
Ruido 75\,$\mu$g/$\sqrt{\text{Hz}}$. Encapsulado LGA-14L de $2.5\times3$\,mm,
masa $\sim$2\,mg. Se monta sobre una flex PCB de poliimida 0.1\,mm adherida
a la piel de la laringe.\\
\\
INMP441 & Micrófono MEMS digital con interfaz I2S. 16\,kHz, 24\,bit.
Conexión directa al periférico SAI del STM32, sin circuitería analógica.\\
\\
TPS7A4700 & LDO de ultra-bajo ruido dedicado exclusivamente a la
alimentación de los sensores. Aísla el ruido de conmutación de la
NUCLEO y la powerbank.\\
\bottomrule
\end{tabularx}

\subsection*{Flex PCB y acople mecánico}

El acelerómetro va montado en una placa flexible de poliimida para minimizar
la carga mecánica sobre la piel:

\begin{itemize}
\item \textbf{Masa total del conjunto:} $\sim$8\,mg (flex + componentes + adhesivo).
\item \textbf{Relación vs.\ piel:} 0.53\,\% de la masa efectiva de la piel laríngea
($\sim$1.5\,g). El sensor sigue fielmente la vibración sin modificarla.
\item \textbf{Adhesivo:} 3M 1526L, acrílico medical grade (ISO 10993),
0.18\,mm de espesor. Refuerzo perimetral con micropore.
\item \textbf{Cableado:} FPC/ZIF pitch 0.5\,mm, 10--15\,cm, con lazo de alivio
de tensión. El cable se fija a la piel 2\,cm antes del conector.
\item \textbf{Frecuencia de resonancia piel--sensor:} $\sim$3.9\,kHz,
fuera de la banda útil para CWT ($<$3\,kHz) y amortiguada por la
viscoelasticidad del tejido.
\end{itemize}

\subsection*{Unidad de procesamiento}

\begin{tabularx}{\textwidth}{l X}
\toprule
\textbf{Periférico} & \textbf{Configuración} \\
\midrule
SPI1 & 10\,MHz, 16-bit, modo 3, DMA circular ping-pong.
Doble buffer de 256 muestras triaxiales.\\
I2S3 & 16\,kHz, 24-bit, DMA circular. Buffer de 128 muestras.\\
TIM2 & Base de tiempo común a 1\,MHz. Ambos flujos DMA comparten
el mismo timer. Cada frame lleva un timestamp en microsegundos.\\
USB OTG\_HS & CDC Virtual COM Port, Bulk endpoint de 512\,bytes.
Streaming sin pérdida de paquetes.\\
\bottomrule
\end{tabularx}

No se ejecuta CWT ni ninguna etapa de inferencia en el microcontrolador.
Todo el procesamiento ocurre en PC durante esta etapa de desarrollo.

\section{Protocolo de comunicación}

Cada trama USB tiene 10\,bytes de cabecera seguidos de \texttt{data\_len} bytes
de payload:

\begin{lstlisting}[style=mini]
Byte 0-1: sync word 0xAA 0x55
Byte 2:   tipo (0x01 = ACC, 0x02 = AUDIO)
Byte 3:   flags
Byte 4-7: timestamp (uint32, microsegundos)
Byte 8-9: data_len (uint16)
Byte 10+: payload
\end{lstlisting}

Los frames se escriben secuencialmente a disco en archivos \texttt{.ssr}.
Cada frame es autodescriptivo: no hay cabecera global, se puede hacer
streaming directo. Los timestamps permiten realinear ambos flujos
a posteriori aunque lleguen en orden arbitrario.

\section{PC tools}

\begin{tabularx}{\textwidth}{l X}
\toprule
\textbf{Herramienta} & \textbf{Qué hace} \\
\midrule
\texttt{ssr-ingest} (Go) & Abre el puerto CDC, detecta sync word 0xAA55,
parsea frames, mantiene un buffer circular de 10\,s en RAM y escribe
asíncronamente a disco. Muestra FPS y tasa de drops en tiempo real.\\
\texttt{ssr-inspect} (Go) & Análisis offline de archivos \texttt{.ssr}:
número de frames, FPS por canal, offset temporal entre flujo ACC y AUDIO.\\
\texttt{ssr-replay} (Go) & Reproduce una sesión grabada, útil para depurar
el pipeline sin el hardware conectado.\\
\texttt{export.py} (Python) & Separa canales, aplica CWT Morlet (50--3000\,Hz,
64 escalas logarítmicas) sobre ACC-X, calcula mel-espectrograma sobre AUDIO,
alinea temporalmente por timestamps y exporta un archivo \texttt{.npz}
listo para consumir en Colab.\\
\bottomrule
\end{tabularx}

\section{Componentes y presupuesto}

\begin{tabularx}{\textwidth}{l l r}
\toprule
\textbf{Componente} & \textbf{Referencia} & \textbf{\$/ud.} \\
\midrule
Dev board STM32 & NUCLEO-H723ZG & 25 \\
Acelerómetro & IIS3DWB (breakout) & 15 \\
Micrófono & INMP441 (breakout) & 2 \\
LDO & TPS7A4700EVM & 6 \\
Cableado + conectores & Dupont + JST SH 1.0mm & 6 \\
Adhesivo médico & 3M 1526L + micropore & 18 \\
Powerbank & 5V 2A & 10 \\
Protoboard + analizador lógico & -- & 13 \\
\midrule
\textbf{Total estimado} & & \textbf{\$85} \\
\bottomrule
\end{tabularx}

\section{Formato del dataset}

Cada sesión de grabación genera:

\begin{lstlisting}[style=mini]
dataset/sessions/{YYYY-MM-DD_HHmmss}/
  metadata.json   --> device, sample_rate, odr, duration, notas
  data.ssr        --> frames binarios secuenciales
  info.md         --> notas de la sesion
\end{lstlisting}

El archivo \texttt{metadata.json} se autogenera con los parámetros de
adquisición. \texttt{export.py} lo convierte en \texttt{.npz} con los
tensores de CWT y mel-espectrograma listos para entrenamiento.

\section{Visión del Producto Final}

SSI evolucionará a un dispositivo completamente autónomo y embebido,
sin dependencia de PC:

\begin{lstlisting}[style=mini]
IIS3DWB ---SPI--- STM32N6 (M55: CWT + NPU: inferencia + M55: sintesis)
                     |
                  MAX98357 --- parlante (voz reconstruida)
\end{lstlisting}

\begin{tabularx}{\textwidth}{l X}
\toprule
\textbf{Característica} & \textbf{Especificación} \\
\midrule
Factor de forma & 50 $\times$ 40 $\times$ 15\,mm, $\sim$32\,g\\
Autonomía & 6\,h inferencia, 400\,días inactivo\\
Procesador & STM32N6 (Cortex-M55 + NPU Neural-Art 1\,TOPS)\\
Sensor & IIS3DWB en flex PCB adherido a laringe\\
Salida & Parlante 30\,mm 2\,W (I2S $\to$ MAX98357)\\
Batería & Li-Po 3.7\,V 500\,mAh, carga USB-C\\
\bottomrule
\end{tabularx}

\subsection*{Mercado objetivo}

El producto está dirigido a:
\begin{itemize}
\item Pacientes con laringectomía parcial o completa
\item Parálisis de cuerdas vocales
\item Afonía funcional o psicógena
\item Comunicación en entornos extremos (bomberos, militares, construcción)
\item Cualquier persona que requiera comunicación silenciosa
\end{itemize}

\section{Roadmap a Muestra TECH 2026}

\begin{tabularx}{\textwidth}{l X}
\toprule
\textbf{Mes} & \textbf{Hito} \\
\midrule
Mayo 2026 & Compra de componentes y montaje en protoboard\\
Junio 2026 & Firmware CMSISv2 funcional (SPI DMA, I2S DMA, USB CDC streaming)\\
Julio 2026 & Primer dataset capturado y export.py validado\\
Agosto 2026 & Entrenamiento de modelo en Colab (variante ReLU, sin atención)\\
Septiembre 2026 & Diseño conceptual PCB producto final + inscripción Muestra TECH\\
15 Octubre 2026 & \textbf{Presentación en Muestra TECH}\\
Post-evento & Fabricación y validación del producto final con STM32N6\\
\bottomrule
\end{tabularx}

\subsection*{Estrategia de evaluación (Muestra TECH)}

\begin{tabularx}{\textwidth}{l c X}
\toprule
\textbf{Criterio} & \textbf{Peso} & \textbf{Enfoque} \\
\midrule
Creatividad e Innovación & 40\% & Uso de acelerómetro ultra-wideband IIS3DWB para captura de vibración laríngea. Producto portátil y discreto para reconstrucción de voz.\\
Nivel Técnico & 25\% & Demo en vivo del prototipo capturando y mostrando stream en PC con FPS estables.\\
Presentación & 20\% & Video del pipeline completo + poster técnico + pitch de 3 minutos.\\
Impacto Social & 15\% & Mercado amplio: pacientes con problemas de voz, entornos ruidosos, comunicación silenciosa. Escalable a múltiples idiomas.\\
\bottomrule
\end{tabularx}

\end{document}
